\section{Setup: symmetry discovery as a nullspace problem}
\label{sec:setup}

\paragraph{Generators and the Lie derivative.} Let $\dot x = F(x)$ be a
dynamical system on $\R^n$ and let $\xi$ be a vector field on $\R^n$. The flow
of $\xi$ is a one-parameter group of state transformations, and it maps
solutions of $\dot x = F(x)$ to solutions if and only if the Lie bracket
\begin{equation}
  \Lc_\xi F \;=\; [\xi, F] \;=\; DF\,\xi \;-\; D\xi\, F
  \label{eq:bracket}
\end{equation}
vanishes identically \citep[Ch.~2]{olver1993applications}. Equation
\eqref{eq:bracket} is the Lie derivative in the sense of
\citet[Eq.~16]{otto2025unified}; the same expression, with $D\xi\,F$ replaced by
the appropriate output action, governs symmetries of a learned map, and with
$D\xi\,F$ deleted it governs invariances of a learned scalar function
(\Cref{sec:invariance}). Everything below applies unchanged in those settings.

We search inside a \emph{candidate class} $\mathcal V =
\mathrm{span}\{\zeta_1, \dots,\zeta_P\}$ of vector fields, writing $\xi_\theta
= \sum_p \theta_p \zeta_p$; in our experiments $\mathcal V$ is the set of
affine generators $\xi(x) = Sx + b$, whose flows are the affine group, or all
polynomial fields of degree at most two. Because \eqref{eq:bracket} is linear
in $\xi$, the exact symmetry algebra inside $\mathcal V$ is the nullspace
$\{\theta : [\xi_\theta,F] \equiv 0\}$, and it is standard to obtain it as the
nullspace of the positive semi-definite Gram operator $\langle \eta,
S\xi\rangle = \int (\Lc_\eta F)^{\!\top} (\Lc_\xi F)\, d\mu$
\citep[\S7]{otto2025unified}.

\paragraph{The target domain.} A symmetry is a statement about a region of
state space. We therefore fix a \emph{target measure} $\nu$, the region over
which the user wants the symmetry to hold, and define all norms as
$\Ltwo{\nu}$ norms. In our experiments $\nu$ is uniform on a box, whose
monomial moments are available in closed form, so \emph{the geometry carries
no Monte-Carlo error}; all statistical uncertainty lives in the model. Keeping
$\nu$ separate from the distribution $\mu$ of the observed states matters:
their mismatch is the incomplete-coverage problem of \Cref{sec:coverage}.

\paragraph{The defect.} Two quantities cancel in \eqref{eq:bracket}. Measuring
the residual relative to their size gives a dimensionless number.

\begin{definition}[Relative equivariance defect]
\label{def:defect}
For a generator $\xi$ and a field $F$,
\begin{equation}
  \defect_\nu(\xi; F) \;=\;
  \frac{\bigl\| \,DF\,\xi - D\xi\,F\, \bigr\|_{\Ltwo{\nu}}}
       {\bigl(\,\|DF\,\xi\|_{\Ltwo{\nu}}^{2}
        + \|D\xi\,F\|_{\Ltwo{\nu}}^{2}\,\bigr)^{1/2}} .
\end{equation}
\end{definition}

\begin{lemma}[Properties]
\label{lem:props}
$\defect_\nu \in [0,\sqrt 2]$ whenever the denominator is nonzero, and
$\defect_\nu(\xi;F)=0$ if and only if $\xi$ generates a symmetry of $F$ on the
support of $\nu$. It is invariant under $\xi \mapsto c\,\xi$ and $F \mapsto
c'F$ for nonzero $c,c'$, hence under any invertible reparameterisation of the
candidate class, and under similarity transformations $x \mapsto cQx$ of the
state ($Q$ orthogonal, $c>0$) with $\nu$ transformed accordingly.
\end{lemma}

The proof is in \Cref{app:proofs}. $\defect$ is the fraction of the bracket's
two terms that fails to cancel, so $\defect = 0.05$ means the two halves of
the bracket agree to $5\%$ in root-mean-square over the target domain. A
generator with $\defect_\nu(\xi;F) \le \delta$ is what we call a
\emph{$\delta$-approximate symmetry}; $\delta$ plays the role that the
singular-value tolerance played before, but now it is comparable across
systems, across candidate bases, and across units.

The denominator vanishes only if $DF\,\xi$ and $D\xi\,F$ both vanish
$\nu$-almost everywhere, which for $\xi \neq 0$ requires a degenerate $F$; the
implementation reports any such direction as uncertified rather than assigning
it a defect.

\begin{remark}[Small eigenvalues are hard to read]
Even with a scale fixed, the smallest eigenvalue of an estimated Gram matrix is
a delicate object: its fluctuations are of a different order from those of the
bulk, and the classical asymptotics \citep{anderson1963asymptotic} and their
random-matrix refinements \citep{johnstone2001distribution,baik2005phase} show
that the map from ``small'' to ``zero in the population'' is not one a fixed
threshold can implement. This is the statistical content of the problem, and it
is why the certificate below is stated for the defect ratio directly rather than
for eigenvalues.
\end{remark}

\begin{remark}[An average, not a supremum]
$\defect_\nu$ is a mean-square quantity, so a generator can have a small defect
while violating the symmetry badly on a small part of the domain, a property it shares with the Gram-operator formulation it refines and with every $L^2$-based criterion. Two responses are available, each with a cost. Choose
$\nu$ to weight the region whose behaviour matters, which is the intended use of
making $\nu$ explicit; or replace the numerator by a supremum norm, which for
polynomials is bounded by a constant times the $L^2$ norm through a Markov-type
inequality, at the cost of that constant. We use the $L^2$ form throughout and
flag this as a modelling choice rather than a technical detail.
\end{remark}

\begin{remark}[Why the normalisation matters]
Without it, the ``small singular value'' criterion is not even well posed:
rescaling one basis generator $\zeta_p$ by $10^{3}$ rescales the corresponding
row and column of the Gram matrix by $10^{3}$ and $10^{6}$, moving eigenvalues
across any fixed threshold. $\defect$ is a functional of the vector field $\xi$
itself and is immune to this.
\end{remark}

\paragraph{Model and data.} We take the standard measured-derivative design of
data-driven model discovery \citep{brunton2016discovering}: pairs
$(x_i, y_i)_{i=1}^N$ with $x_i \sim \mu$ and
\begin{equation}
  y_i \;=\; F(x_i) + \varepsilon_i , \qquad \varepsilon_i \sim N(0, \sigma^2 I_n)
  \ \text{i.i.d.},
  \label{eq:model}
\end{equation}
and a model class $F_\beta(x) = \sum_{i,a} \beta_{i,a}\, x^{a} e_i$ of
polynomial fields of degree at most $d$, so that $\beta \in \R^Q$ with
$Q = n\binom{n+d}{d}$. Ordinary least squares gives $\bhat$ and, in this
homoskedastic Gaussian model, an \emph{exact} finite-sample confidence
ellipsoid; \Cref{sec:robustness} replaces it with a sandwich form and reports
what changes. Throughout, $\sigma$ is reported relative to the
root-mean-square size of $F$ over the target domain, so ``noise level $0.1$''
means the noise standard deviation is a tenth of the typical size of the
right-hand side.

\begin{lemma}[Bilinear structure]
\label{lem:bilinear}
For polynomial $F_\beta$ and polynomial candidate class, there are matrices
$L_\theta, M^{(1)}_\theta, M^{(2)}_\theta \in \R^{K\times Q}$, linear in
$\theta$ and independent of the data, such that
\begin{equation}
  \defect_\nu(\xi_\theta; F_\beta)^2
  \;=\; \frac{\|L_\theta \beta\|^2}
             {\|M^{(1)}_\theta \beta\|^2 + \|M^{(2)}_\theta \beta\|^2},
  \qquad L_\theta = M^{(1)}_\theta - M^{(2)}_\theta ,
  \label{eq:ratio}
\end{equation}
where $\|\cdot\|$ is the Euclidean norm and $K$ is the dimension of the space of
possible brackets. Consequently
$\min_{\theta \ne 0} \defect_\nu(\xi_\theta;F_\beta)^2$ is the smallest
generalised eigenvalue of the pencil $(C(\beta), D(\beta))$ with
$C = L^{\!\top}L$ and $D = M^{(1)\top}M^{(1)} + M^{(2)\top}M^{(2)}$ assembled
over the basis of $\mathcal V$.
\end{lemma}

The matrices are obtained by expanding \eqref{eq:bracket} in a monomial basis
and absorbing the $\Ltwo{\nu}$ moments into a symmetric square root; the
construction is exact linear algebra and is given in \Cref{app:construction}.
The reduction to a generalised eigenproblem is itself a small improvement on
the standard recipe: the plug-in symmetry directions become independent of how
the candidate basis is scaled.
